\chapter{Introduction}\label{demo-framing}
A sensor can only report the water that reaches it. In a branching river, this simple fact creates a placement problem: a downstream instrument may observe many sources but conceal where a disturbance began. This chapter defines the imagined problem without borrowing a field site or a real monitoring record.
\section{Research questions and objectives}\label{demo-framing-question}
We ask how a small instrument budget should be distributed across branches when disturbances travel downstream and observations can fail. The target is a transparent decision procedure, not a prediction for an existing watershed.
\subsection{Scope of the tracer model}\label{demo-framing-scope}
The experiment concerns a dissolved tracer with an arbitrary unit of concentration. No toxicity threshold is implied. All candidate sites are assumed accessible and equally costly unless a scenario explicitly changes that assumption.
\subsection{Event detection criteria}\label{demo-framing-evidence}
An event is detected when a simulated concentration exceeds a fixed threshold at an available sensor. We record event coverage separately from arrival delay so that late detection cannot masquerade as prompt warning.
\section{Representing downstream transport through branching networks with shared observation sites}\label{demo-framing-network}
Edges point downstream. A node may receive two upstream flows, but no edge returns water to an ancestor. This excludes tides and recirculating reservoirs by construction.
\begin{figure}[htbp]
\centering
\NetworkDiagram{West branch}{East branch}{Outlet}
\caption{An invented two-branch river, with a confluence and one downstream observation site.}\label{demo-fig-network}
\end{figure}
\subsection{Synthetic measurement units}\label{demo-framing-units}
Distance is measured in synthetic reaches and time in ticks. A reach has no fixed conversion to kilometres; a tick has no fixed conversion to hours. These choices prevent unwarranted geographic interpretation.
\begin{table}[htbp]
\centering
\begin{tabular}{lll}
\hline
Symbol & Meaning & Illustrative range \\
\hline
$n$ & Candidate sites & 7--31 \\
$b$ & Sensor budget & 1--5 \\
$t$ & Arrival tick & 0--24 \\
\hline
\end{tabular}
\caption{Symbols used throughout the fictional experiment.}\label{demo-tab-symbols}
\end{table}
\section{Scope and limitations}\label{demo-framing-boundaries}
An efficient arrangement on a graph is not necessarily a feasible deployment. Power, access permissions, maintenance, and habitat disturbance remain outside the simulated objective.
\subsection{Consistency checks}\label{demo-framing-checks}
The first check is conceptual: every numerical statement must identify the simulated world that produced it. Where the thesis uses a hand-worked example rather than a computed trial, that distinction is made explicit.
